\rsection{Logik \& Entscheiden}
\slide{Begrifflichkeiten}{
    \begin{itemize}
        \item \b{Entscheiden:} Auswahl einer Handlungsalternative unter mehreren, um ein Ziel zu erreichen (zielorientiert, rational).
        \item \b{Wahrheitswert:} Der logische Status einer Aussage (wahr oder falsch); Elementaraussagen haben einen eindeutigen Wahrheitswert.
        \item \b{Junktoren/Konnektoren:} Logische Verknüpfungen (z. B. UND, ODER, WENN-DANN), die Elementaraussagen zu komplexen Aussagen verbinden.
        \item \b{Utility Function (Nutzenfunktion):}
        \begin{itemize}
            \item Formalisierung rationaler Entscheidungen (Homo Oeconomicus).
            \item Berechnung: Summe aus (Gewichtung eines Merkmals $\times$ Ausprägung des Merkmals) für jede Alternative
            \item Ziel: Maximierung des subjektiven Nutzens.
        \end{itemize}
    \end{itemize}
}
\slide{Rationalitäts-Konzepte}{
    \begin{itemize}
        \item \b{Unbounded Rationality (Rationalität 1):}
        \begin{itemize}
            \item Streng regelbasiertes, logisches System.
            \item Ignoriert kognitive Beschränkungen; Ziel ist die optimale Lösung (Maximierung).
            \item Entspricht eher Computern oder ökonomischen Modellen.
        \end{itemize}
        \item \b{Bounded Rationality (Rationalität 2):}
        \begin{itemize}
            \item Assoziatives, heuristisches System.
            \item Berücksichtigt begrenzte Ressourcen (Zeit, Wissen, Kapazität); Ziel ist eine zufriedenstellende Lösung.
            \item Entspricht menschlichem Verhalten im Alltag.
        \end{itemize}
        \item \b{Satisficing:} Eine Heuristik (Mischung aus sufficing und satisfying), bei der die erste Option gewählt wird, die ein Anspruchsniveau erfüllt ("gut genug"), statt nach der optimalen Lösung zu suchen.
    \end{itemize}
}
\slide{Arten des logischen Schließens}{
    \begin{enumerate}
        \item Deduktion: Schluss vom Allgemeinen auf das Besondere (Regel $\rightarrow$ Fall). Wenn die Prämissen wahr sind, ist die Konklusion zwingend wahr (sicherer Schluss).
        \item Induktion: Schluss vom Besonderen auf das Allgemeine (Beispiele $\rightarrow$ Regel). Die Konklusion ist wahrscheinlich, aber nicht sicher (unsicherer Schluss).
        \item Abduktion: Schluss von einem Sachverhalt auf eine mögliche Ursache. Es ist ein Vorschlag zur Erklärung, aber sehr unsicher.
    \end{enumerate}
}
\slide{Aussagenlogik \& Schlussregeln}{
    \begin{itemize}
        \item \b{Modus Ponens:} Wenn $P \rightarrow Q$ und $P$ gegeben ist, folgt $Q$ (logisch einfach, Fehlerquote gering).
        \item \b{Modus Tollens:} Wenn $P \rightarrow Q$ und $\text{nicht-}Q$ gegeben ist, folgt $\text{nicht-}P$ (logisch schwieriger, Fehlerquote hoch).
        \item \b{Wason's Task (Wahlaufgabe):}
        \begin{itemize}
            \item Experiment zur Überprüfung von bedingten Schlussfolgerungen (Karten mit Vokalen/Zahlen).
            \item Befund: Menschen tun sich schwer mit dem Modus Tollens bei abstrakten Inhalten (drehen oft die falschen Karten um, Bestätigungsfehler).
            \item Cheater-Detection: Wird die Aufgabe in einen sozialen Kontext eingebettet (Betrüger entlarven, z. B. Alkoholkonsum unter 18), steigt die Lösungsquote drastisch an (evolutionär verankertes Modul?).
        \end{itemize}
    \end{itemize}
}
\slide{Theorien \& Heuristiken}{
    \begin{itemize}
        \item \b{Mentale Logik (Beweistheorie):} Theorie, dass Menschen über ein inneres Set an logischen Regeln verfügen und diese direkt auf Prämissen anwenden.
        \item \b{Mentale Modelle:} Theorie, dass Menschen konkrete Vorstellungen (Modelle) des Problemraums bauen und diese inspizieren/variieren, statt abstrakte Regeln anzuwenden.
        \item \b{Heuristiken:}
        \begin{itemize}
            \item Availability Heuristik: Die Wahrscheinlichkeit eines Ereignisses wird danach eingeschätzt, wie leicht Beispiele aus dem Gedächtnis abgerufen werden können (Verfügbarkeit).
            \item Representativeness Heuristik: Wahrscheinlichkeitseinschätzung basiert auf der Ähnlichkeit zu einem Prototypen, oft unter Missachtung von Basisraten.
            \item Conjunction Fallacy: Fehlannahme, dass eine spezifische Kombination von Eigenschaften (Bankangestellte \& Feministin) wahrscheinlicher ist als eine einzelne Eigenschaft (Bankangestellte), ausgelöst durch Repräsentativität.
            \item Base Rate Fallacy: Das Ignorieren der Grundwahrscheinlichkeit (Basisrate) in einer Population zugunsten spezifischer Informationen.
        \end{itemize}
    \end{itemize}
}